% test-025.tex -- comandos \spProm, \spmedia, \spmediana, \spmoda, \spdesv y \spvarz
%
% Compilar:  lualatex-dev test-025.tex
% Validar:   show-pdf-tags --xml test-025.pdf | rnv-wrapp
%            verapdf --flavour ua2 --format text test-025.pdf
%
% Cada caso es un parrafo numerado, para poder referirse a el al
% escuchar el PDF. Los comentarios "doc:" indican la lectura que
% documenta el manual de usuario; en el resto, la lectura se revisa
% escuchando. Las entradas invalidas (que detienen la compilacion)
% no van aqui.
\DocumentMetadata
  {
    lang=es-CL,
    pdfversion=2.0,
    pdfstandard=ua-2,
    tagging=on,
    tagging-setup={math/setup={mathml-SE}},
  }
\documentclass{article}
\usepackage[spanish]{babel}
\usepackage{unicode-math}
\usepackage{spintent}
\usepackage{hyperref}
\hypersetup
  {
    pdfdisplaydoctitle = true,
    pdftitle           = {test-025: estadistica},
  }
\pagestyle{empty}
\begin{document}

\section*{A. El comando spProm}

% doc: «promedio de x»
1. Una letra: $\spProm{x}$.

% doc: «promedio de y»
2. Otra letra: $\spProm{y}$.

% doc: «promedio de x sub uno»
3. Una letra con subíndice: $\spProm{x_{1}}$.

% doc: «promedio de x más y»
4. Una expresión: $\spProm{x+y}$.

% doc: «media de x»
5. Con read-arg: $\spProm[read-arg={media de}]{x}$.

% doc: sin lectura
6. Con silent: $\spProm[silent=true]{x}$.

\section*{B. Los símbolos}

% doc: «media aritmética»
7. Media aritmética: $\spmedia$.

% doc: «mediana»
8. Mediana: $\spmediana$.

% doc: «moda»
9. Moda: $\spmoda$.

% doc: «desviación estándar»
10. Desviación estándar: $\spdesv$.

% doc: «varianza»
11. Varianza: $\spvarz$.

\section*{C. Llave sigma}

% doc: «desviación estándar»
12. Con sigma en spdesv: $\spdesv[sigma=true]$.

% doc: «varianza»
13. Con sigma en spvarz: $\spvarz[sigma=true]$.

\section*{D. Llaves sub-index y read-sub}

% doc: «media aritmética de x»
14. Con sub-index: $\spmedia[sub-index={x}]$.

% doc: «media aritmética sub x»
15. Con sub-index y read-sub: $\spmedia[sub-index={x},read-sub=true]$.

% doc: «mediana de x»
16. Mediana con sub-index: $\spmediana[sub-index={x}]$.

% doc: «moda de x»
17. Moda con sub-index: $\spmoda[sub-index={x}]$.

% doc: «desviación estándar de x»
18. Desviación estándar con sub-index: $\spdesv[sub-index={x}]$.

% doc: «varianza de x»
19. Varianza con sub-index y sigma: $\spvarz[sub-index={x},sigma=true]$.

\section*{E. Llaves read-arg y silent}

% doc: «desviación estándar poblacional»
20. Con read-arg: $\spdesv[read-arg={desviación estándar poblacional}]$.

% doc: «varianza muestral»
21. Con read-arg y sub-index (el subíndice no se lee): $\spvarz[read-arg={varianza muestral},sub-index={x}]$.

% doc: sin lectura
22. Con silent: $\spmedia[silent=true]$.

\section*{F. Dentro de fórmulas}

% doc: «promedio de x es igual a x sub uno más x sub dos más x sub tres, partido por tres»
23. Promedio: $\spProm{x} = \dfrac{x_{1}+x_{2}+x_{3}}{3}$.

% doc: «desviación estándar es igual a la raíz cuadrada de varianza»
24. Desviación estándar y varianza: $\spdesv = \sqrt{\spvarz}$.

% doc: «media aritmética es igual a promedio de x»
25. Media aritmética y promedio: $\spmedia = \spProm{x}$.

% doc: «media aritmética, mediana y moda»
26. Los tres juntos: $\spmedia$, $\spmediana$ y $\spmoda$.

\end{document}
